\chapter{اللقاء الخامس: وجوه التأويل وبداية التفسير}

\section{الوجوه التي يُوصل بها إلى معرفة تأويل القرآن}
السلام عليكم ورحمة الله وبركاته. الحمد لله رب العالمين، والصلاة والسلام على نبينا محمد وعلى آله وصحبه وسلم. هذا هو الدرس الثالث من اجتماعنا على دراسة كتاب الإمام \rawi{الطبري} رحمه الله «جامع البيان عن تأويل آي القرآن».

وصلنا بحمد الله إلى فصل من أهم فصول المقدمة، وهو: القول في الوجوه التي من قِبَلِها يوصل إلى معرفة تأويل القرآن. 

\isnad{قال أبو جعفر رحمه الله:} «ونحن قائلون في البيان عن وجوه مطالب تأويله. قال الله جل ثناؤه لنبيه محمد صلى الله عليه وسلم: \begin{ayah}وَأَنزَلْنَا إِلَيْكَ الذِّكْرَ لِتُبَيِّنَ لِلنَّاسِ مَا نُزِّلَ إِلَيْهِمْ وَلَعَلَّهُمْ يَتَفَكَّرُونَ\end{ayah}... وقد تبين ببيان الله جل ذكره أن مما أنزل الله من القرآن على نبيه ما لم يُوصَل إلى علم تأويله إلا ببيان الرسول صلى الله عليه وسلم... وأن منه ما لا يعلم تأويله إلا الله الواحد القهار، وذلك ما فيه من الخبر عن آجال حادثة وأوقات آتية كوقت قيام الساعة... وأن منه ما يعلم تأويله كل ذي علم باللسان الذي نزل به القرآن، وذلك إقامة إعرابه ومعرفة المسميات بأسمائها».

\begin{fawaid}[وجوه مطالب التأويل عند الطبري]
لخص \rawi{الطبري} الوجوه التي يُفهم بها القرآن إلى ثلاثة أوجه رئيسية:
\begin{enumerate}
    \item \textbf{ما لا يُعلم إلا ببيان النبي صلى الله عليه وسلم:} وهذا يستلزم عدم القول في كتاب الله في تفاصيل الأحكام والشرائع إلا بعد العلم بالسنة النبوية.
    \item \textbf{ما استأثر الله بعلمه:} كوقت قيام الساعة ومجيء أشراطها الكبرى، وهذا لا يُطلب تأويله لأنه لا سبيل للوصول إليه.
    \item \textbf{ما يُعلم من لسان العرب:} وهو معرفة معاني الألفاظ وتصاريف الكلام، وهذا يوجب العناية الفائقة بلغة العرب لفهم القرآن.
\end{enumerate}
\end{fawaid}

\separator

\section{القول في القرآن بالرأي وتوجيه آثار السلف}
\isnad{ساق الطبري رحمه الله عدة أخبار في النهي عن تفسير القرآن بالرأي، منها:}
\isnad{عن} \rawi{ابن عباس} \isnad{أن النبي صلى الله عليه وسلم قال:} \begin{hadith}من قال في القرآن برأيه فليتبوأ مقعده من النار\end{hadith}. \isnad{وعن} \rawi{أبي بكر الصديق} \isnad{أنه قال:} «أي أرض تقلني وأي سماء تظلني إذا قلت في كتاب الله عز وجل برأيي أو بما لا أعلم».

\isnad{قال أبو جعفر:} «وهذه الأخبار شاهدة لنا على صحة ما قلنا، من أن ما كان من تأويل القرآن الذي لا يُدرك علمه إلا بنص بيان رسول الله... فغير جائز لأحد القيل فيه برأيه، بل القائل في ذلك برأيه وإن أصاب عين الحق فيه فمخطئ في فعله... لأن إصابته ليست إصابة موقن، وإنما هي إصابة خارص وظان، والقائل في دين الله بالظن قائل على الله ما لا يعلم».

\begin{fawaid}[قاعدة في الخطأ في التفسير]
القائل في كتاب الله بغير علم، حتى وإن وافق قوله الصواب، فهو مخطئ آثم لفعله ما نُهي عنه من التقول على الله بغير علم. وهذا يتوافق مع ما قرره \rawi{الشافعي} من أن من تكلف ما جهل فموافقته للصواب غير محمودة، وخطؤه غير معذور فيه.
\end{fawaid}

\subsection{توجيه آثار السلف في الإحجام عن التفسير}
ثم ساق \rawi{الطبري} أخباراً عن بعض التابعين (مثل \rawi{سعيد بن المسيب} و\rawi{عبيدة السلماني} و\rawi{الشعبي}) تُفيد إحجامهم عن التفسير.

وقد وجه \rawi{الطبري} صنيعهم توجيهاً دقيقاً؛ فقال إن إحجامهم لم يكن إنكاراً لأن يكون للقرآن تفسير، بل كان «إحجام خائف ألا يبلغ باجتهاده ما كلفه الله من إصابة الصواب». فقاس إحجامهم عن التفسير على إحجام بعض الأئمة عن الفتيا في النوازل تورعاً وخوفاً من الخطأ، لا تحريماً لأصل الاجتهاد والفتيا.

\separator

\section{شروط المفسر وأدواته}
بعد أن بيّن \rawi{الطبري} وجوه التأويل، وضع خلاصة جليلة تُمثل (أصول التفسير) وتحدد من هو الذي يحق له أن يتكلم في القرآن.

\isnad{قال رحمه الله:} «فأحق المفسرين بإصابة الحق في تأويل القرآن... أوضحهم حجة فيما تأول وفسر... مما كان تأويله إلى رسول الله صلى الله عليه وسلم، من أخبار رسول الله الثابتة عنه (إما بنقل مستفيض أو نقل عدول)... أوضحهم برهاناً فيما ترجم وبين مما كان مُدرَكاً علمه من جهة اللسان (بالشواهد من أشعارهم واللغات المستفيضة)... بعد أن لا يكون خارجاً في تأويله عما تأول وفسر السلف من الصحابة والأئمة من التابعين».

\begin{fawaid}[آلة المفسر ومضمار التفسير]
حدد \rawi{الطبري} آلة المفسر في ثلاثة أصول:
\begin{itemize}
    \item النقل الثابت عن النبي صلى الله عليه وسلم.
    \item الحجة والبرهان من لغة العرب وشواهد أشعارهم.
    \item أن يلتزم بمضمار الإجماع، فلا يخرج بتفسيره عن مجموع أقوال السلف من الصحابة والتابعين (بأن يحدث قولاً يباين أقوالهم).
\end{itemize}
\end{fawaid}

\separator

\section{أسماء القرآن وسوره}
انتقل \rawi{الطبري} لبيان أسماء القرآن وأسماء سوره، فذكر أن الله سماه: (قرآناً، وفرقاناً، وكتاباً، وذكراً).
ورجح أن كلمة (قرآن) هي مصدر من القراءة (بمعنى المقروء)، كالكفران والخسران، ورجح هذا على قول \rawi{قتادة} بأنه من الضم والجمع.

\begin{fawaid}[منهج الطبري في تأخير المسائل لمظانها اللائقة]
أثناء كلامه عن أسماء السور (كالطوال، والمئين، والمثاني، والمفصل)، أشار إلى اختلاف العلماء في معنى (المثاني). ولكنه لم يُفصل المسألة هنا، بل قال: «وسنذكر أسماء قائلي ذلك... عند انتهائنا إلى قوله تعالى: \textit{ولقد آتيناك سبعاً من المثاني}». وهذه من أبرز ميزات منهجه: أنه يترك التفصيل في المسألة إلى الموضع القرآني الأليق بها ولا يقحمها في المقدمات.
\end{fawaid}

\subsection{معنى السورة والآية}
بيّن \rawi{الطبري} أن السورة في لغة العرب مأخوذة من «المَنْزِلة من منازل الارتفاع» كسُور المدينة، فكأن سور القرآن منازل من الشرف مرتفعة. واستدل بقول \rawi{النابغة الذبياني}:
\begin{quote}
ألم ترَ أن الله أعطاك سورةً *** ترى كُلَّ مَلْكٍ دونها يتذبذبُ
\end{quote}

أما (الآية)، فرجح أنها إما بمعنى (العلامة) التي يُعرف بها تمام ما قبلها، وإما بمعنى (القصة) التي تتلو القصة.

\separator

\section{الشروع في تفسير الفاتحة والاستعاذة}
ثم بدأ الإمام \rawi{الطبري} في التفسير العملي، مبتدئاً بأسماء سورة الفاتحة (أم القرآن، فاتحة الكتاب، السبع المثاني). وبيّن وجه تسميتها بـ «أم القرآن» لأن العرب تُسمي كُل جامعٍ لأمرٍ (أُمّاً)، كأم الرأس وأم القرى.

\subsection{تأويل الاستعاذة والشيطان والرجيم}
\isnad{قال:} «الاستعاذة الاستجارة، وتأويل \textit{أعوذ بالله من الشيطان} أستجير بالله دون غيره. والشيطان في كلام العرب: كل متمرد من الجن والإنس والدواب، سُمي بذلك لمفارقة أخلاقه أفعال سائر جنسه... والرجيم هو فعيل بمعنى مفعول (مرجوم)، أي الملعون المشتوم المطرود».

\subsection{تأويل البسملة والفرق بين الرحمن والرحيم}
بيّن \rawi{الطبري} أن البسملة تتضمن فعلاً محذوفاً يُقدر بحسب الفعل الذي يُشرع فيه، فكأنه يقول عند القراءة: «أقرأ مبتدئاً بتسمية الله».

ثم عالج مسألة الفرق بين (الرحمن) و(الرحيم)، مؤصلاً أنه لا توجد كلمة في كتاب الله تُغني عن الأخرى أو تقوم مقامها من كل وجه.

\begin{fawaid}[الفرق بين الرحمن والرحيم]
رجح \rawi{الطبري} أن (الرحمن) بناء يدل على عموم الرحمة لجميع خلقه (المؤمن والكافر في الدنيا)، وأن (الرحيم) يدل على خصوص الرحمة ببعض خلقه (بالمؤمنين في الدنيا والآخرة).
\end{fawaid}

\begin{fawaid}[قاعدة في تقديم الاسم على الصفة]
أجاب \rawi{الطبري} عن سؤال: لماذا قُدِّم اسم (الله) ثم (الرحمن) ثم (الرحيم)؟ بأن من شأن العرب إذا أرادوا الخبر عن شيء أن يُقدموا الاسم العَلَم أولاً ثم يُتبعوه بالصفات والنعوت. ولأن اسم (الله) خاص لا يتسمى به غيره مُطلقاً، قُدّم، ثم تلاه (الرحمن) الذي مُنع التسمي به أيضاً، ثم (الرحيم) الذي يجوز وصف بعض المخلوقين به.
\end{fawaid}

بهذا نختم مقدمة الكتاب، ونشرع في الدروس القادمة في تفسير الآيات تفصيلاً بإذن الله.